\section{Exhaustion of the capacity conditions}
\label{sec:capacity}

The balance and band constructions cover complementary shapes. A small
sum of the first four rows or columns gives balance support. Otherwise
both orientations have room for the path in the band criterion. Failure
of that criterion then bounds both row and column prefixes, forcing an
upper bound on the whole diagram's area. We use this contradiction for
large parameters and for most intermediate parameters. Nineteen pairs
remain, for which we enumerate targets using the proved support tests.

Throughout this section, $M\ge9$ and $r>0$, and $M,r,b,s,L,K$ have
the meanings fixed above. Put
\[
 c=2b+2,\qquad T=K+r,
 \qquad H_i=\sum_{a=1}^i h_a,\qquad W_j=\sum_{a=1}^j w_a,
\]
where $h$ and $w$ are the column heights and row lengths of a target
partition of $n$, extended by zeros. Since $L$ is self-conjugate, we
may apply every test to either orientation of the target.

By Proposition~\ref{prop:balance-four}, any target not already
supplied satisfies
\begin{equation}
 H_4\ge K,\qquad W_4\ge K.
 \label{eq:capacity-four-assumption}
\end{equation}

\subsection{Two opposite prefixes}

\begin{lemma}
\label{lem:capacity-rectangle}
Let $i,j\ge1$. If $H_i\le U$ and $W_j\le V$, then
\begin{equation}
 n\le\max\left(U+V-1,
              \left\lfloor\frac{UV}{ij}\right\rfloor\right).
 \label{eq:capacity-rectangle}
\end{equation}
\end{lemma}

\begin{proof}
If the cell in row $j$, column $i$ is absent, no cell lies outside
both the first $i$ columns and the first $j$ rows. Those two strips
cover the diagram and share its first cell, giving $n\le U+V-1$.
Otherwise their intersection is an $i$-by-$j$ rectangle. Put
$x=h_i\ge j$ and $y=w_j\ge i$. Every cell outside both strips has
row index at most $x$ and column index at most $y$, so the remaining
corner has at most $(x-j)(y-i)$ cells. Since $ix\le U$ and $jy\le V$,
\[
 n\le U+V-ij+(x-j)(y-i)
 \le U+V-ij+(U/i-j)(V/j-i)=\frac{UV}{ij}.
\]
Integrality gives the asserted floor.
\end{proof}

\begin{lemma}
\label{lem:capacity-large}
If $M\ge22$, a target satisfying
Equation~\eqref{eq:capacity-four-assumption} passes the band test in
one of its two orientations.
\end{lemma}

\begin{proof}
The choice $d=4$, $q=8$ is admissible in
Proposition~\ref{prop:band-test}, because
$dq+7+d=43\le2M-1$. If the band test fails in both orientations,
then
\[
 \sum_{i=1}^8\max\left(0,\left\lfloor\frac{h_i-4}{2}\right\rfloor\right)
 \le r-1,
\]
and likewise for $w$. The termwise inequality
\begin{equation}
 x\le2\max\left(0,\left\lfloor\frac{x-f}{2}\right\rfloor\right)+f+1
 \quad(x\in\mathbb Z_{\ge0})
 \label{eq:capacity-floor-bound}
\end{equation}
therefore gives $H_8,W_8\le S:=2r+38$.
Lemma~\ref{lem:capacity-rectangle} implies
$n\le\max(2S,S^2/64)$. On the other hand,
\[
 n-2S\ge\frac{M^2-5M-160}{2}>0,
 \qquad
 64n-S^2\ge23M^2-28M-1508>0.
\]
For the first inequality substitute $r\le3M/2+2$; for the second
write $64n-S^2=32M(M+1)-4r^2-24r-1444$ and make the same
substitution. Both displayed polynomials are positive at $M=22$
and increase thereafter. This is a contradiction.
\end{proof}

For the remaining intermediate cases, retain the least adequate
core height in each orientation. Let $d$ be least with $H_d\ge K$
and $e$ least with $W_e\ge K$; by
Equation~\eqref{eq:capacity-four-assumption}, both belong to
$\{1,2,3,4\}$. Define
\[
 q_f=\min\left(8,\left\lfloor\frac{2M-8-f}{f}\right\rfloor\right)
 \quad(1\le f\le4),
\]
and the following finite set of pairs consisting of a prefix length
and an upper bound for its sum:
\begin{equation}
 \begin{split}
 \mathcal D_{M,r}(d,e)
 ={}&\{(d-1,2M-2):d>1\}\\
 &{}\cup\{(q_f,2r-2+q_f(f+1)):e\le f\le4,\ q_f>0\}.
 \end{split}
 \label{eq:capacity-D}
\end{equation}
If neither band test succeeds, every pair $(i,U)$ in this set
satisfies $H_i\le U$. Indeed minimality of $d$ gives its first
element, and $W_f\ge K$ for $f\ge e$; failure of the band test with
this $f,q_f$, followed by Equation~\eqref{eq:capacity-floor-bound},
gives the other elements. Similarly,
$\mathcal D_{M,r}(e,d)$ bounds the row prefixes.
The indices are crossed because enough room in the first $f$ rows
allows the band criterion to test the excess heights of columns,
and conversely after transposition.

\begin{lemma}
\label{lem:capacity-intermediate}
The first-four balance test and the band tests cover every target
when $9\le M\le21$, $r>0$, except possibly for
\begin{equation}
 (M,r)\in
 \{(10,8),\ldots,(10,17)\}
 \cup\{(12,14),\ldots,(12,20)\}
 \cup\{(14,22),(14,23)\}.
 \label{eq:capacity-exceptions}
\end{equation}
\end{lemma}

\begin{proof}
Use $1\le r\le\lfloor M/2\rfloor$ for odd $M$ and
$1\le r\le3M/2+2$ for even $M$. After removing
Equation~\eqref{eq:capacity-exceptions}, there are 177 pairs $(M,r)$.
For each pair and each of the 16 choices of $d,e$, the following
integer inequality holds:
\begin{equation}
 \min_{\substack{(i,U)\in\mathcal D_{M,r}(d,e)\\
                 (j,V)\in\mathcal D_{M,r}(e,d)}}
 \max\left(U+V-1,\left\lfloor\frac{UV}{ij}\right\rfloor\right)<n.
 \label{eq:capacity-intermediate-check}
\end{equation}
This is a bounded integer calculation involving 2,832 assertions.
The function \texttt{intermediate} in
Appendix~\ref{app:capacity-code} evaluates exactly
Equation~\eqref{eq:capacity-intermediate-check}; its complete execution
verified every assertion. If a target failed the first-four and band
tests, its least $d,e$ would give both the upper bounds in
Equation~\eqref{eq:capacity-D}. Equation~\eqref{eq:capacity-intermediate-check}
would then contradict Lemma~\ref{lem:capacity-rectangle}.
\end{proof}

\subsection{The nineteen exceptional cases}

For the remaining parameters we first assemble induced supports of full
degree. When these do not settle a target, we allocate cells to a
constituent with at most four rows on the big pair and test the residual diagram
against the lower components. A recursive enumeration checks that every
target is covered by one of these supports or by the band criterion.
All arithmetic is exact; the tests certify the sufficient support
conditions, rather than approximate Kronecker coefficients.

For a multiset $A$ of positive integers, write
\[
 |A|=\text{its number of entries},\qquad
 \|A\|=\sum_{a\in A}a,\qquad
 P_A(k)=\sum_{a\in A}\min(k,a).
\]
By Lemma~\ref{lem:young-support}, inducing trivial representations of
the sizes in $A$ has precisely the support of partitions
$\gamma\vdash\|A\|$ satisfying
\begin{equation}
 \sum_{j=1}^k\gamma_j^t\le P_A(k)\quad(k\ge1).
 \label{eq:capacity-young}
\end{equation}
We call this the support associated with $A$. For integers $u,j>0$,
let $\operatorname{pack}(u,j)$ be the multiset of $j$ integers
$\lfloor(u+i)/j\rfloor$, $0\le i<j$, with zero entries omitted.
This is the balanced division from Section~\ref{sec:balance}.
Every use below has $u\ge j$, so the code's fixed-length tuples
have no zero entries. The implementation in
Appendix~\ref{app:capacity-code} follows these constructions with
integer lists and sets.

For a consecutive increasing list $J$ of middle lengths, let
$\mathcal P(J)$ be the following collection of multisets. If $|J|$
is even, pair consecutive entries $u,u+1$, replacing each pair by
$\operatorname{pack}(2u+1,4)$. If $|J|$ is odd, choose an entry with
an even number of predecessors, retain it as a single entry, and
perform this pairing separately on the two remaining even lists.
The empty list gives the empty multiset. Thus every member of
$\mathcal P(J)$ has the support supplied by those middle components,
by Lemma~\ref{bal:neighboring} and
Lemma~\ref{lem:balance-separation}. The function \texttt{bal}
enumerates this collection.

\paragraph{Full-degree supports.}
Construct a collection $\mathcal B$ as follows. For each row of
Table~\ref{tab:capacity-full-supports} having the prescribed value of
$s$, and each $R\in\mathcal P(J)$, take the multiset union of $R$
with the displayed parts. The indices $(t,x)$ match those in the code.
Omit a row if $x>c$.

\begin{table}[ht]
\centering
\small
\setlength{\tabcolsep}{5pt}
\begin{tabular}{@{}cccl@{}}
\hline
$s$ & $(t,x)$ & Parts adjoined to $R$ & $J$\\
\hline
$1$ & $(0,0)$ & $\operatorname{pack}(T,4)\uplus(c,1)$
    & $(3,\ldots,M-2)$\\
$1$ & $(0,1)$ & $\operatorname{pack}(T,4)\uplus\operatorname{pack}(c+1,2)$
    & $(3,\ldots,M-2)$\\
$2$ & $(2,0)$ & $\operatorname{pack}(T+2,4)\uplus(c)$
    & $(3,\ldots,M-2)$\\
$2$ & $(2,3)$ & $\operatorname{pack}(T+2,4)\uplus\operatorname{pack}(c+3,2)$
    & $(4,\ldots,M-2)$\\
\hline
\end{tabular}
\caption{Full-degree support multisets. The symbol $\uplus$ denotes
multiset union; the row with $x=3$ uses the length-three block in
the short component and removes it from $J$.}
\label{tab:capacity-full-supports}
\end{table}

Here $t=2$ combines the two singleton cells with the big pair;
$x=1$ combines the singleton with the length-two blocks; and
$x=3$ combines the length-three block with those blocks.
The singleton left in the first row contributes the part $1$.
Every $A\in\mathcal B$ has $\|A\|=n$, and its support is
supplied by Lemmas~\ref{bal:neighboring},~\ref{bal:two-paths},
and~\ref{bal:odd-strip}, combined by
Lemma~\ref{lem:balance-separation}.

\paragraph{Supports below the big pair.}
These full-degree supports sometimes settle the target immediately.
For the remaining targets we seek a columnwise allocation
\[
 h_i=\delta_i+(h_i-\delta_i),\qquad
 0\le\delta_i\le\min(4,h_i),\qquad \sum_i\delta_i=T,
\]
with at least one odd $\delta_i$ when $s=2$.
The big pair supplies the partition with column-height multiset
$(\delta_i)$ by Lemmas~\ref{bal:neighboring} and~\ref{bal:two-paths}.
We therefore need lower-component support for the residual partition,
whose column heights are the decreasing rearrangement of
$(h_i-\delta_i)$ and whose size is $n-T$.
Lemma~\ref{lem:lr-addition} will combine these two constituents.

Next construct a collection $\mathcal C$ for the components other
than the big pair. Write $u=c+3+s$, and let $\mathcal G$ consist of
the partitions of $u$ with at most three rows obtained from a
partition of $c$ with even row lengths by successively adding
horizontal strips of sizes $3$ and $s$.
Proposition~\ref{prop:three-row} supplies all of $\mathcal G$ on the
blocks of lengths at most three. For partitions of equal size,
write $\alpha\unrhd\beta$ for dominance of row lengths. Retain a
partition $a$ of $u$ with at most three rows if
\begin{equation}
 \alpha\unrhd a\quad\Longrightarrow\quad\alpha\in\mathcal G
 \quad\text{for every $\alpha\vdash u$ with at most three rows}.
 \label{eq:capacity-low-mask}
\end{equation}
No partition of more than three rows can dominate $a$, since its
first three rows have sum strictly less than $u$. Thus
Lemma~\ref{lem:young-support} says that the entire support associated
with the row lengths of $a$ is available on the lower blocks.
For each retained $a$ and each
$R\in\mathcal P((4,5,\ldots,M-2))$, put their multiset union in
$\mathcal C$. Every resulting multiset has total $n-T$ and supplied
support. We may delete $A$ from $\mathcal C$ if its decreasing
rearrangement strictly dominates that of another member $B$:
every partition dominating $A$ also dominates $B$, so this deletion
does not shrink the union of their supports. Deleting all such
members simultaneously is harmless: every descending chain in this
finite strict dominance order terminates at a retained member whose
support contains the supports of its predecessors.

These definitions explain the construction of \texttt{B} and
\texttt{C} in the code. The function \texttt{parts} lists all
partitions with at most three parts. For such partitions $a,b$, the
function \texttt{strip} tests the interlacing inequalities
\[
 a_1\ge b_1\ge a_2\ge b_2\ge a_3\ge b_3\ge0,
\]
which are exactly the horizontal-strip condition. The calls have
the prescribed size differences $3$ and $s$. The list \texttt{bad}
is the complement of $\mathcal G$, so the test excluding any
$\alpha\in\texttt{bad}$ dominating $a$ is exactly
Equation~\eqref{eq:capacity-low-mask}. The final dominance deletion
is the valid deletion just described.

\paragraph{Safe pruning of partial diagrams.}
The available supports have now been specified. The next bound determines
when an initial column list already certifies every possible completion.

\begin{lemma}
\label{lem:capacity-prefix-pruning}
Let $A$ be a nonempty multiset of positive integers with total $N$, and let
$a=(a_1,\ldots,a_l)$ be a nonempty prefix of the column-height list
of a partition of $N$. If
\begin{equation}
 \min\left(N,\sum_{v=1}^{\min(i,l)}a_v+(i-l)_+a_l\right)
 \le P_A(i)\qquad(1\le i<\max A),
 \label{eq:capacity-prefix-bound}
\end{equation}
then every completion of $a$ has the support associated with $A$.
For a complete column-height list followed by a zero, these inequalities
are equivalent to that support condition.
\end{lemma}

\begin{proof}
Every unknown entry is at most $a_l$. The left side of
Equation~\eqref{eq:capacity-prefix-bound} is therefore an upper
bound for the corresponding prefix sum of any completion. For
$i\ge\max A$, the right side is already $N$, so the omitted
inequalities hold automatically. If the prefix includes a final
zero, it has no nonzero extension and the upper bounds are exact.
\end{proof}

The predicate \texttt{bound(A,a)} is exactly
Equation~\eqref{eq:capacity-prefix-bound}; it returns false on an
empty prefix. The predicate \texttt{F(a,b)} applies the band test,
using known column and row prefixes and the largest admissible
$q=\min(8,\lfloor(K-7-d)/d\rfloor)$ for each $d=1,\ldots,4$.
The lists contain exact initial row and column lengths. Summing
only the known entries therefore gives lower bounds for both sums
in the band criterion, since every capacity summand is nonnegative.
Extending either prefix cannot invalidate a successful test.
Thus the predicate \texttt{stop(a,b)}, which accepts either band
orientation or any support from $\mathcal B$, certifies all
completions of these prefixes. Its use before the diagram is fully
constructed is justified by Lemma~\ref{lem:capacity-prefix-pruning}.

\paragraph{Certifying a completed allocation.}
If prefix pruning does not settle a target, the next routine allocates
cells to the big pair and tests the residual diagram against the lower
supports. Only successful allocations are used.

\begin{lemma}
\label{lem:capacity-hit}
If the code returns \texttt{hit(h)=True} on a complete column-height
list followed by zero, the target is supplied by the big pair and
the other balance components.
\end{lemma}

\begin{proof}
The routine constructs allocations $\delta_i$ from the original
columns, with $0\le\delta_i\le\min(4,h_i)$. If $s=1$, the protected
index is the final zero, to which it allocates zero. If $s=2$, it
chooses a column and fixes its allocation to $z=1$ or $3$; no later
step changes this column. Every other step decreases the residual
of an eligible column by one, stopping at residual
$\max(0,h_i-4)$. Entering the loop's \texttt{else} clause means
exactly $T=K+r$ cells have been allocated. Taking the positive
allocations as column heights gives a partition with at most four rows
and size $T$,
and when $s=2$ it has the protected odd-height column.
Its support on the big pair is supplied
by Lemma~\ref{bal:neighboring} when $s=1$ and
Lemma~\ref{bal:two-paths} when $s=2$.

The residual columns are nonnegative and have total $n-T$, the
total of every member of $\mathcal C$. A successful \texttt{bound}
call on their decreasing rearrangement, which retains a final zero,
certifies an available residual constituent by
Lemma~\ref{lem:capacity-prefix-pruning}. Indeed, the original appended
zero is protected when $s=1$; when $s=2$, it is neither eligible for
the positive allocation $z$ nor for a later decrement. Sorting does
not remove that zero. For each original column separately,
the height-$h_i$ column occurs in the induction of the columns of
heights $\delta_i$ and $h_i-\delta_i$. Applying
Lemma~\ref{lem:lr-addition} to these single-column containments
supplies the original target. No optimality claim about the greedy
choice of decrements is needed: only its successful allocations are
used.
\end{proof}

\paragraph{Exhausting the targets.}
We have justified what each successful test certifies. It remains to
show that the recursive search visits every target not already certified.

\begin{lemma}
\label{lem:capacity-enumeration}
For each pair in Equation~\eqref{eq:capacity-exceptions}, the
procedure \texttt{gen} exhausts all targets, up to transposition and
the already justified pruning conditions.
\end{lemma}

\begin{proof}
At depth $k$, the lists $a,b$ give the first $k$ full row and column
lengths, and \texttt{left} is the area after removal of these
successive diagonal hooks. The remaining southeast partition, if
nonempty, has a first row of length $x$ and a first column of length
$y$. Necessarily
\[
 1\le x\le\texttt{left},\quad
 1\le y\le\texttt{left}+1-x,\quad
 \texttt{left}\le xy.
\]
If $k>0$, decreasing full row and column lengths further require
$x\le a_k-k$ and $y\le b_k-k$. These are exactly the loop bounds:
the rectangle inequality is written
$y\ge\lceil\texttt{left}/x\rceil$. The next lists append $x+k$
and $y+k$, and the remaining area becomes
$\texttt{left}+1-x-y$, subtracting the next hook.
At the root the additional bound $y\le x$ selects the orientation
whose first row is at least its number of rows. Every target has
such an orientation, and all final tests allow both orientations.
Every actual partition therefore follows a permitted path.

Conversely, the row and column constraints make the corresponding
Frobenius arm and leg lists strictly decreasing and nonnegative.
Every completed path consequently specifies a partition. At
\texttt{left=0}, its rows beyond the first $k$ are
\[
 \#\{j\le k:b_j\ge i\},\qquad k<i\le b_1,
\]
and its columns are given by the symmetric formula. These are
exactly the two calls to \texttt{fill}; the appended zero is for
the support tests. Hook sizes account for the entire area $n$.

The recursion terminates since a nonempty hook removes at least
one cell. A successful \texttt{stop} discards only certified
completions. The other early rejection, at $k\ge4$, is used only
when a first-four sum is below $K$; those four entries are then
fixed and Proposition~\ref{prop:balance-four} applies. At each remaining
leaf the program verifies
$\texttt{hit(u) or hit(v)}$ for the complete conjugate lists.
By Lemma~\ref{lem:capacity-hit}, a successful assertion certifies
that target as well.
\end{proof}

The calls \texttt{run(m,r)} in Appendix~\ref{app:capacity-code}
were executed for every pair in
Equation~\eqref{eq:capacity-exceptions}. All nineteen calls
completed with every assertion satisfied; none was omitted using
the optional earlier rectangle test. This verifies the remaining
finite assertion in Lemma~\ref{lem:capacity-enumeration}.
As supplementary checks on the implementation, the accompanying
\texttt{capacity\_independent.py} compared diagonal-hook enumeration
with ordinary partition enumeration for every $1\le n\le40$
(113,696 emitted partitions in the selected orientation), checked
the three-row Pieri support sets in all nineteen cases by a separate
columnwise test, checked the totals of the multisets in $\mathcal B$
and $\mathcal C$, and checked 748,716 prefix-bound instances for
$1\le n\le16$.
These checks all passed; the preceding invariants, rather than
their smaller degree ranges, justify exhaustive coverage in the
nineteen required cases.

\begin{proposition}
\label{prop:capacity-exhaustion}
For every admissible $M\ge9$ and $r>0$, every target partition of
$n=N_M+2r$ is supplied by the band and balance constructions for
$L$.
\end{proposition}

\begin{proof}
Proposition~\ref{prop:balance-four} handles a small first-four sum in
either orientation. For the remaining targets,
Lemma~\ref{lem:capacity-large} handles $M\ge22$, and
Lemma~\ref{lem:capacity-intermediate} handles $9\le M\le21$
outside Equation~\eqref{eq:capacity-exceptions}.
The nineteen verified calls, justified by
Lemmas~\ref{lem:capacity-hit} and~\ref{lem:capacity-enumeration},
handle exactly those remaining pairs.
\end{proof}
